

\begin{table}[h!]
\footnotesize
\centering
\caption{Customer Statements and Interpreted Needs for Thermostat Design}
\label{tab:interpreted_needs}
\begin{tabular}{|p{2.5cm}|p{5cm}|p{5cm}|p{2.5cm}|}
\hline
\textbf{Question/Prompt} & \textbf{Customer Statement} & \textbf{Interpreted Need} & \textbf{Variable Type} \\
\hline
Typical uses & I have to manually turn it on and off when it gets too hot or cold. & The thermostat maintains a comfortable temperature without requiring user action. & Functional \\
\cline{2-4}
& Each time I want to change the temperature, I need to adjust both thermostats in the house. & Any user input need not be made in multiple locations. (I) & Functional \\
\hline
Likes—current model & I like that I can change the temperature if the setting is too high. & The temperature setting is easy to control manually. & Usability \\
\cline{2-4}
& It didn’t cost a fortune. & The thermostat is affordable to purchase. & Cost \\
\hline
Dislikes—current model & I’m too lazy to learn how to figure it out. & The thermostat requires little or no user instruction or learning. & Usability \\
\cline{2-4}
& I sometimes forget to turn it off when we leave the house. & The thermostat saves energy when no one is home. (†) & Energy Saving \\
\cline{2-4}
& Sometimes the buttons don’t register a push and I have to press it repeatedly. & The thermostat definitively registers any user inputs. & Reliability \\
\hline
Suggested improvements & I would like my iPhone to adjust my thermostat. & The thermostat can be controlled remotely without requiring a special device. & Connectivity \\
\cline{2-4}
& I like to shift quickly between different options like energy saving or ultra comfort. & The thermostat responds immediately to differences in occupant preferences. & Responsiveness \\
\hline
\end{tabular}
\end{table}